\section{Related Work}
\label{sec:related-work}

\subsection{Distributed Trace Representations}

The effectiveness of learning-based trace anomaly detection depends
strongly on how distributed traces are represented before analysis.
Distributed traces are variable-length execution records with service
invocation paths, parent-child relationships, and heterogeneous latency
behavior. Prior work has therefore explored several representations for
trace analysis. TraceAnomaly represents each trace as a fixed-dimensional
Service Trace Vector (STV), where features correspond to service call
paths and values encode latency-related behavior~\cite{liu2020traceanomaly}.
Graph-based approaches instead preserve more of the trace topology.
DeepTraLog constructs a trace event graph that combines invocation
structure with log events~\cite{zhang2022deeptralog}, while TraceVAE
models microservice traces using a graph variational autoencoder
designed for tree-structured trace data~\cite{xie2023tracevae}. Other
recent approaches incorporate additional telemetry such as profiling
metrics or resource utilization to enrich trace representations
\cite{panahandeh2024serviceanomaly,alomari2025contextaware}.

These representations involve different trade-offs. Fixed-dimensional
vectors such as STVs are lightweight and directly interpretable because
each feature maps to a concrete service path, but they may collapse
structural information. Graph-based and multi-modal models preserve more
execution structure and operational context, but typically require more
complex learning architectures and may reduce transparency of individual
feature contributions. Our work builds on the interpretability and
simplicity of STVs, but studies a different question: whether anomaly
scores derived from raw STV features are systematically influenced by
trace structural complexity.

\subsection{Trace-Based Anomaly Detection and Localization}

Early approaches to microservice anomaly detection often relied on
manually defined rules or service-specific heuristics to identify
abnormal executions~\cite{liu2020traceanomaly}. As microservice systems
increased in scale and complexity, these rule-based techniques became
harder to maintain, motivating learning-based methods that model normal
trace behavior directly from data. TraceAnomaly learns the distribution
of normal STVs using service-level deep Bayesian networks and assigns
higher anomaly scores to traces that deviate from the learned normal
behavior~\cite{liu2020traceanomaly}. Because STV features correspond to
service paths, this representation also provides service-path-level
diagnostic information.

Subsequent methods improve trace anomaly detection by using richer
representations or more expressive models. DeepTraLog combines trace and
log information in a graph neural architecture with a deep support vector
data description objective~\cite{zhang2022deeptralog}. TraceGra uses
graph deep learning to model trace-based anomaly detection in
microservice systems~\cite{chen2023tracegra}. ServiceAnomaly combines
distributed traces with profiling metrics in an annotated graph to
detect abnormal microservice executions~\cite{panahandeh2024serviceanomaly}.
Recent work has also explored context-aware trace graph embeddings and
resource-metric autoencoders for anomaly detection and localization
\cite{alomari2025contextaware,xing2025multidim}.

These methods primarily focus on improving anomaly detection accuracy,
capturing a richer trace structure, or enhancing fault localization. In
contrast, our work focuses on the behavior of the anomaly scoring process
itself. Specifically, we examine whether STV-based anomaly scores can be
biased by structural properties of traces, such as the number of spans or
services, even when the traces are operationally normal.

\subsection{Research Gap: Bias in STV-Based Anomaly Scoring}

Existing trace anomaly detection studies have advanced the state of the
art through improved trace representations, graph neural architectures,
deep generative models, and multi-modal telemetry fusion. However, they
typically evaluate detection accuracy or localization quality, while the
relationship between anomaly score and trace structural complexity is
rarely examined directly. This leaves open an important question for
STV-based anomaly scoring: does a high anomaly score reflect abnormal
service-path behavior, or can it also reflect the fact that a trace is
structurally larger than other traces?

To the best of our knowledge, prior STV-based trace anomaly detection
studies have not explicitly quantified whether raw STV anomaly scores
are correlated with trace-complexity measures such as span count and
service count. Our work addresses this gap by measuring the dependency
directly, proposing a lightweight residual scoring method designed to
reduce it, and evaluating the resulting effect on both trace-complexity
bias and controlled latency-anomaly scoring performance.